% RESEARCH_DRAFT_SHA256: ce0811089e10bda8815de0f21b4749c8db16ddd966048e95f2203394bb5d368a
% Layout and source-reference wrappers expanded for the web renderer.
\documentclass[11pt]{article}
\usepackage[margin=25mm]{geometry}
\usepackage{fontspec,unicode-math}
\setmainfont{Latin Modern Roman}
\setmathfont{Latin Modern Math}
\usepackage{amsmath,amssymb,mathtools,xurl,hyperref}
\hypersetup{colorlinks=true,urlcolor=blue,linkcolor=blue}
\setcounter{secnumdepth}{0}
\setlength{\parindent}{0pt}
\setlength{\parskip}{5pt}
\setlength{\emergencystretch}{3em}
\newcommand{\sref}[2]{\hyperlink{source:#1:#2}{[S#2]}}
\newcommand{\cataloguescope}[1]{\paragraph{Source record.}#1}
\newcommand{\floor}[1]{\left\lfloor#1\right\rfloor}
\begin{document}
\section*{Research-attempt provenance}
This standalone research attempt retains the complete original definition body below,
including its historical source and preparation notes. Those notes describe the
earlier catalogue preparation; the new attempt and its verification are identified
separately after that body. No original ID, catalogue status, or source statement is changed.

% UnsolvedMath ID 20001164; source catalogue code AIM-COMPUTATION-0002
\section{A malformed composite record and an Artin--Schreier certificate family}\label{20001164}
{\small UnsolvedMath ID 20001164 · Computer Science}
\subsection{Definitions and mathematical statement}

\paragraph{Shared notation.}Standard number systems are $\mathbb N,\mathbb Z,\mathbb Q,\mathbb R,\mathbb C$. Specialized notation and index ranges follow the cited source and the definitions below.


A polynomial over a field is irreducible if it has no factorization into two positive-degree polynomials. An irreducibility certificate is auxiliary data checked by a verifier to establish this property. Sparse encoding lists only nonzero monomials and coefficients. A finite field $\mathbb F_q$ has $q=p^n$ elements; sums involving powers $x^{p^i}$ are linearized expressions, while products of such powers can represent quadratic forms over $\mathbb F_p$.
\paragraph{Record scope and qualifications.}
The catalogue title is retained as metadata. Named certificates, contracts, bounds or frameworks in that title are not asserted as results of the original source extract.
The source text joins distinct questions and loses exponent, summation, and probability notation. Its claimed hardness and fraction of inputs cannot be interpreted as a precise algorithmic promise from this extraction. The supplied catalogue marks this record open; its present research status was not independently checked. The extraction is labelled unrecoverable and is retained as a literal source record; no exact conjecture is reconstructed.
The source text below is rendered literally where necessary; damaged formulas and composite extracts are not silently repaired into a new mathematical claim.
The supplied catalogue marks this record as unrecoverable. Its extract and terminology are documented, but a complete original problem has not been reconstructed.
\paragraph{Statement recorded in the supplied source.}
1. Finding better certificates for irreducibility (a) Are there certificates for irreducibility that one can verify faster than the existing irreducibility tests? (Victor Miller's Question) (b) Can you exhibit families of sparse polynomials for which you can find shorter certificates? 2. Find a polynomial-time algorithm in log( q), where q = pn, that solves

∑n−1

> i,j =0

ai,j xpi +pj

+ ∑n−1

> i=0

bixpi

+ c in Fq (or shows no solutions) and works for 1

> poly

of the inputs, and never lies. (We know that this is NP-complete.) (Jintai Ding's Question) 13 Factoring over Q[x]

3.1 Open
Question \href{https://aimath.org/WWN/polyfactor/polyfactor.pdf}{[S2]}
\paragraph{Source status.}
The supplied catalogue labels this source record open; that metadata was checked against the supplied snapshot on 2026-10-11. The label is preserved for provenance and is not an independent verification of present open status or mathematical correctness.
\subsection{Short English statement}
The extraction combines irreducibility-certificate questions and a finite-field equation-solving question.
\subsection{Sources}
\begin{itemize}
\item[S1] Ulam AI and the UnsolvedMath Contributors, UnsolvedMath: A Curated Collection of Open Mathematics Problems, record 20001164 (AIM-COMPUTATION-0002). Catalogue curation and metadata: CC BY 4.0. The source extract is retained; definitions, English summaries and qualifications are editorial additions.
\url{https://huggingface.co/datasets/ulamai/UnsolvedMath}.
\item[S2] AIM Workshop Problem Lists, source linked by catalogue record AIM-COMPUTATION-0002.
\url{https://aimath.org/WWN/polyfactor/polyfactor.pdf}.
\end{itemize}
\paragraph{Preparation and review.}
The original source and UnsolvedMath contributors retain attribution for the source extract. Codex (GPT-6; exact serving revision not exposed) assisted with definitions, English summaries, scope notes and formatting. No new human mathematical authorship or independently verified solution is asserted. Permanent source ID 20001164 and code AIM-COMPUTATION-0002 checked against the complete supplied registry and current repository. Distinct upstream records, including repeated questions and broad or damaged extracts, are retained for complete coverage of the source CS category. A repeated record is not a new independent conjecture. Complete mathematical-equivalence and current-literature review has not been performed. Preparation and final review follow contribution prompts C and D, with source gaps disclosed. Structural and rendering checks do not establish a new theorem or open status.
\label{end:20001164}

\section{Literature and context check}
The AIM polynomial-factorization PDF was consulted at Section 2.2. It
contains separate certificate and finite-field equation questions. The damaged
combined extraction and its unspecified success fraction remain unrepaired as
an original target. The following is a clearly declared sparse-family test,
not a reconstruction of the missing promise or a solution of the composite.
\section{Research attempt: A sparse-family generalization fails under Frobenius iteration}
The supplied prior preparation treats the familiar family $T^p-T-a$ using
trace. An attempted enlargement replaces $p$ by $p^k$ while keeping a three-term
description, hoping that a nonzero base-field trace would still certify
irreducibility. It already fails over the prime field.

Let $a\in\mathbb F_p^\times$ and $f(T)=T^{p^k}-T-a$. If $\alpha$ is any root
in an algebraic closure, then
$\alpha^{p^k}=\alpha+a$. Inductively,
\[
 \alpha^{p^{jk}}=\alpha+ja\qquad(j\ge1),
\]
since Frobenius fixes $a$; at $j=p$ this yields
$\alpha^{p^{pk}}=\alpha$. Thus its orbit under the ordinary $p$-Frobenius
has length at most $pk$. To see directly that its minimal polynomial has degree
at most this length, multiply $T-\alpha^{p^i}$ over one complete orbit. Frobenius
permutes the factors and fixes the coefficients; an algebraic-closure element
fixed by $c\mapsto c^p$ lies in $\mathbb F_p$, since $T^p-T$ has exactly its
$p$ prime-field roots. The orbit product is therefore a polynomial over
$\mathbb F_p$ annihilating $\alpha$. This proves the degree bound without
assuming irreducibility of $f$.

Consequently if $p^k>pk$, $f$ cannot be irreducible: any root has a nonconstant
minimal polynomial of smaller degree dividing $f$ (division with remainder
and evaluation proves divisibility). This is a systematic obstruction, not a
single accidental factor. In the first example $p=2,k=3,a=1$,
\[
 T^8+T+1=(T^2+T+1)(T^6+T^5+T^3+T^2+1)
 \quad\hbox{over }\mathbb F_2.
\]
The trace of $a$ to the prime field is $1$, yet the polynomial is reducible.
Multiplying the displayed factors verifies every coefficient.

For $k=1$ the degree obstruction does not apply; it says $p\le p$. For
$a=0$ the polynomial has the evident root zero, a separate boundary case.
Therefore this calculation rejects only the proposed larger sparse family,
not the known prime-degree trace criterion. Short sparse encoding alone does
not preserve a valid irreducibility certificate as exponents grow.

The other bundled equation-solving question is untouched: a factorization
obstruction for this family provides neither a promised fraction of solvable
instances nor a never-incorrect polynomial-time solver for quadratic
linearized expressions over the extension field.
\section{Verification and adversarial audit}
Frobenius iteration was checked for all integers $j$ symbolically, and the
factorization was multiplied by Python over $\mathbb F_2$, giving low-to-high
coefficients $(1,1,0,0,0,0,0,0,1)$. The proof distinguishes prime field
coefficients from arbitrary extension-field coefficients, where $a$ need not
be fixed by one Frobenius step.
\paragraph{Audit corrections.} A tempting extension of the prior Artin--Schreier calculation was explicitly falsified rather than asserted. Both the damaged extraction and untouched second question remain disclosed.
\paragraph{External computation.} Python 3.14.2 exact binary-coefficient multiplication of the two displayed factors produced T^8+T+1.
\section{Scope and remaining gaps}
The record is composite and unrecoverable as a single precise conjecture.
No new certificate-size or verification-complexity theorem for the intended
general family is proved, and the equation-solving promise remains unspecified.
The first gap in this attempted sparse-family route is an actually irreducible
family with a sound cheap verifier under a fully charged field representation.
\paragraph{Integration review.} Codex (GPT-6; exact serving revision not exposed) reviewed the written arguments and expanded layout and source-reference wrappers. Repository metadata and the production web renderer were checked using Python 3.13 and JavaScript with MathJax 3.2.2. This model-assisted review is not independent human proof verification; no PDF compilation or formal proof checking is asserted.

\section{Completion Estimate}
\noindent\textbf{COMPLETION: 0\%}
The calculations below the preserved source statement investigate a concrete route,
but establish no substantiated new progress toward the entire original target.
Known special cases and failed sufficient conditions are not credited as a solution.
\section{References}
The original source references and catalogue attribution are preserved in the
Sources section above. The following URLs were actually consulted for this attempt;
retention of an additional original reference does not claim fresh consultation of it.
\begin{itemize}
\item \url{https://aimath.org/WWN/polyfactor/polyfactor.pdf} (consulted 2026-10-11). 
\end{itemize}
\section{Model and contributor disclosure}
OpenAI gpt-6-astra, requested reasoning configuration ultra, generated this research
attempt on 2026-10-11. The exact serving revision is not exposed. Codex GPT-6
performed orchestration and repository-format review; its exact serving revision
is also not exposed. The mathematical route, boundary checks and final adversarial
self-audit followed contribution prompts A then B in
\texttt{docs/contribution\_prompts.md}, after reading \texttt{CONTRIBUTING.md}.
Python 3.14.2 assembled the standalone documents; any mathematical computation is
specified in Verification. No theorem prover was used. No human mathematical
authorship, substantial human proof correction, or independent human proof checking
is asserted. Model-assisted review, source retrieval, and structural checks do not
establish correctness or novelty. Full private reasoning transcripts are unavailable
as reproducibility artifacts; the complete arguments, inputs and outcomes actually
used are stated here or in the companion batch audit.
\section{Final status}
Unresolved. This is a research attempt with an explicit failed route or restricted
calculation, not a full solution, not a current-status certification, and not a claim
that the elementary material is new. The catalogue's original source-status label
is preserved independently of this attempt status.
\end{document}
